% Split from 3_sylvia.tex by cut-sheet v2/v3 — content below is the source scene, header adjusted
\chapter{The Pain of our Penumbras}
% \section*{An Arms race}
Northbound. Piccadilly line. Ernest is on the Tube, northbound, still carrying Sylvia’s voice like a faint electrical hum in his skull. The Airbnb on Caledonian Road is waiting for him with its low ceiling, damp towel smell, a mattress that remembers other men. He doesn’t mind. He has slept in worse places. The assorted opposite have just tiredness in common.

\begin{figure}[H]
\includegraphics[width=\linewidth]{Illustrations/vign-tube.png}
\caption{Ernest's view of on the way home}
\end{figure}

The train exhales into King’s Cross. Doors open and the tube refills itself. Two women step in together. Not demonstrative but close enough that their alignment is obvious. That they are orbiting a common center. Ernest registers their coupling without curiosity. Adjusts. Slides one seat down so they can sit side by side.

He stands half a second too quickly, gestures lightly, and shifts one seat down so they can sit together. Always the corrective movement. Always the ethical adjustment.

Which places him directly beside a man with arm all ready. Already. Not resting. Occupying. His forearm feels corded: all wiry strength rather than bulk. 

The geometry is tight. He could yield. There is room to concede just an inch.

This man on his left has been there since Leicester Square. Ernest registers him in the tunnel reflection only in fragments: shorn head, close-cropped beard, a stillness that feels less calm than withheld. Head, shaven, cleaved. His face looks almost erased. No softness. No slackness. 

Ernest sits tight arm on his arm. 

The train pulls out. The tunnel darkens the windows to mirrors.

In the reflection, he studies the man a little further. Calculates. Frame, shoulder width, probable center of gravity. The wiry ones are dangerous. Tendon and leverage. Quick. But not too large. He reckons, that is if it came to it, that he could take him.  It’s strange how the mind performs this arithmetic without invitation.

He thus finds himself resting his left forearm along a taught arm already there. Not aggressive. But very present.

Ah. The muscle beneath him hardens further. Now it is taking form.

Opening his phone he scrolls everything and nothing. The screen glows uselessly. He reads nothing while feeling everything.Arm on arm. This is now the axis of the universe.

At Russell Square the carriage compresses further. A lateral lurch throws weight sideways and presses him more firmly into the contact. The voice arrives low and so very close.

“If you don’t fucking get your arm off mine I’ll end you, you cunt.”

The carriage inhales.

Ernest does not move but feels his heart answer not with fear exactly, but voltage.

The woman he made space for, urgently ushers him: “Just go sit over there.” Escape offered. Civility disguised as prudence. But to move now would rewrite the geometry. No. Not facing saving. To sit opposite would be concession disguised as civility.

He turns only his head. Slowly. Owl-like. Deliberately. He meets the man’s eyes.
They are angry, yes. But they are not abyssal. Not feral. There is theatre in them. Assertion more than annihilation. The rest of him still, heart racing and face neutral Ernest determines this guy is not unhinged. All performance over psychosis. In the reflection, the shaved head looks like something machined. But up close the little pores betray him. He is man. Tense. Wired.

Ernest glances down at their shared territory with a small shrug as if puzzled by the laws of physics. As if to say: we both know what is happening here. He refuses to show the recoil coiled inside him.

Stops come. Bodies shift. Each jolt a further press on the arm beneath his. Each jolt drives them fractionally closer to tumult. Pressure becomes message. And so the man finally stands.

Ernest thinking: this is it. But the man walks away timidly but angrily offering up,

“You need a lesson, you little cunt.”

He steps off. Doors close. Air returns to the carriage.

“Oh what a horrible man,” someone says.

Ernest looks at his phone as if resuming something important. Breath measured. Suitably, apparently Unruffled. He who has survived car accidents. Hospitals. Sylvia’s gaze. The collapse of whole cosmologies.

And yet .... what is unsettling him is not the threat but his readiness to react. How quickly his body answered. How instantly his old circuitry lit up. Sylvia had only just restored his nihilism of clean abstraction; that nothing matters, that meaning is projection. But an arm on an arm had so just mattered. Pressure mattered. Territory and dominance mattered. He had calculated bone density and leverage from a reflection of a tunnel window. For a man who claims nothing matters, that is a revealing reflex.

The train pulls north and the platform slips away. But when the train enters the tunnel and the windows turn to mirrors again, he studies the carriage behind him in reflection.
Is he still there? No. But reflections are unreliable. Movement overlays movement. Faces slide into each other. The mind fills what the eye cannot fix.

Caledonian Road. He stands. Steps off. The platform feels wider than usual. He does not rush. That would concede something to the universe. He walks at a measured pace. Spine upright. But he listens. Footsteps behind him? There are always footsteps behind you in a station. He does not turn immediately for turning too quickly is an admission.
Halfway up the stairs he glances down ever so casually, a side swipe of a look, as if checking a message. No shaved head. Just commuters. A student. A woman with a suitcase.

Still. The thought enters fully formed: What if he did not get off? What if this was not about armrest geometry at all, but about selecting a mark? The body replays the arithmetic. Distance to exit. Escalator bottleneck. Angle of approach. Where would he strike first? It does not arrive as a full memory. It arrives as force. The lurch of metal. The sideways compression. That specific weight of impact when physics decides for you. It all comes back. His ribs tightening under the seatbelt. The glass fracture pattern like frozen lightning. Sylvia’s name not yet spoken but already in consequence. he then retraces the Tube’s braking shudder just moments ago and they become indistinguishable from that earlier deceleration. For a flicker, the station tiles blur into a roadside verge. He inhales sharply.

This is ridiculous. No one is following him. Even if the man were, what then? A confrontation at the ticket barrier? A theatrical shove? A lesson administered in fluorescent light? The mind loves to complete patterns. Arm on arm becomes pursuit. Pursuit becomes violence. Violence becomes destiny.

Ernest's pace has meant now he has gained on a woman ahead. He notices her without wanting to. Early thirties. Dreadlocks tied back. Then the rest of her. He registers it, files it, feels the small, uninvited tightening that comes with the noticing.

Now he notices an old man approaching from the other direction. Head slightly dipped. A schnauzer at the end of a loose lead. Small thing.  Already displeased with the world.

Ernest’s mind once locked onto elbows drops fully into the present.

The dog goes for her. A burst of barking. Sharp. Excessive. Directed. The old man does not stop. Does not apologise. Walks on as if nothing has happened.

And then the thought ,he realises he wants the dog to bark at him as well. He feels the pull of it. A need for things to be even, even if they are going to be unpleasant.

Please. Don’t single her out. Then, more bluntly. Please don’t be a racist dog. I cannot be dealing with a racist dog tonight. Not after those contained angers. The dog turns. Looks at him.
A moment. Measuring. Then it barks. Properly. At him. Ernest feels something of a release.

Yes. Good. That will do. He keeps walking. Over the small rise.

He walks toward the damp Airbnb, all narrow brick frontage, a door that sticks slightly before yielding. Halfway there he has the sudden conviction that someone is pacing him on the opposite side of the road. He does not look. If he looks, the narrative hardens. He keeps walking. Heart rate elevated but contained. All of it, he tells himself, is projection. Just as Sylvia argued. Meaning is an overlay and threat is an interpretation. The body reacts to shadows and then builds theology around them. The lesson, the pursuit, the future confrontations are all mind. Like everything.

He reaches the door. Unlocks it. Enters the narrow hallway that smells faintly of detergent and old carpet. Locks the door behind him. And only then does he allow himself a single glance through the frosted glass. Nothing. Of course nothing. He stands there longer than necessary. What disturbs him is not the man. It is how quickly the world becomes adversarial in his head. How easily geometry becomes narrative. Sylvia had restored his nihilism : that clean claim that nothing has inherent meaning. But the body, the nervous system refuses nihilism. Insisting rather on pattern. 

Ernest turns the key in the Caledonian Road door.At the same time, across the city, Sylvia is being raised carefully from the passenger seat of Hector’s car.
Hector is competent in the way of men who have made a life out of competence. The chair folds. The brace is adjusted. The door opens with a careful, practiced weight.

“Bit of traffic tonight,” he says lightly, as they arrive in Holland Park.
He always narrates the obvious. Hector talks. About parking restrictions. About a client. About a small inefficiency in the lift maintenance at her building that he has “had words about.” His voice is steady, mildly managerial. Protective.

She feels it grate. Because he is kind. Because he is constant. Because he contains no risk. Her mind drifts back, unwillingly to Ernest perhaps on a tube platform, yes a platform.

Something dark flickers and the revenge thought returns without announcement.  Let him be visited, she thinks. Let something come at him that he cannot out-think. Let there be a force that does not respect his geometry, his detachment. Let him feel the lack of control he has metered out so coolly.

The thought is sharp. Bitter. Almost satisfying.

Hector continues talking. “ ... and I told them, if you’re going to resurface the courtyard you might as well do it properly ...”

She hears herself exhale. There is guilt immediately. Not theatrical guilt. Just a tightening. Ernest did not engineer the accident. Physics did that. Youth did that. Momentum did that. Her wish for his visitation, though feels … disproportionate.

Cruel. The trees arch darkly above them. Lights in the apartment blocks glow warm and arranged.

“I’ll walk you in,” Hector says. Of course he will.  The perpetual readiness. The way he manages the world into safe increments.

“Honestly, Hector, I can manage,” she snaps ... sharper than intended.

He pauses. Only for a second. “Of course,” he says gently. But he continues unfolding the chair. The irritation is misdirected and she knows it even as she feels it. He is not the architect of her confinement, he is the custodian of it.

The doors of the building open. The lift hums. The apartment smells faintly of lavender and radiator heat. Inside, she removes her gloves slowly. The revenge fantasy dissolves under its own ugliness. What she wanted, ever so briefly, was not Ernest harmed. It was Ernest shaken. Ernest forced out of abstraction. Ernest subject to something he could not master.

And the uglier truth beneath that: she wants him to feel the contingency she lives with daily. The helpless geometry of her altered bones.

Hector is adjusting the cushions behind her without being asked and she closes her eyes for a moment.

“Thank you,” she says quietly. But she can already feel the faint resentment coiling again, not at him, but at the equilibrium of it all.

The apartment settles. Radiator ticks. Distant traffic hums below Holland Park. The city is muted here, filtered through glass and height. Sylvia wheels herself to the window and looks down at the clean geometry of the streetlamps. Circles of sodium light. Even spacing. Predictable intervals. Order imposed on darkness.

She turns from the window and moves to her desk. Papers are spread there, her notes from work with John. Diagrams half-inked. Arcs, intersections, invariants. The cross-ratio sketched in one margin. Four points projected along a line. The constancy under transformation. She runs her fingers lightly over the page. There are quantities that remain fixed even when perspective shifts. There are relations preserved under distortion. That is what she works on now with John, in extracting invariants from systems that appear chaotic under naïve coordinates.

She closes her eyes. The accident was not teleology. It was trajectory intersecting trajectory. Her bitterness is an attempt to solve an equation with anger as a coefficient that destabilizes everything.

She breathes heavily. In her work with John, they had been refining a projection argument; of how structures survive transformation if you identify the right invariants.
She whispers the word to herself: Invariant. Her injury is not erased by that word.
Ernest is not redeemed by it. But the system becomes clearer.If he were undone tomorrow, it would not restore symmetry. The revenge fantasy collapses under arithmetic.

She turns a page and begins annotating again. Line AB projected to A'B'. Ratio preserved.
She feels the anger thinning, replaced not by forgiveness, but by proportion. Ernest is a finite perturbation in a much larger space. Her work persists. Her mind persists.
Her collaboration with John, so steady, rigorous, unsentimental and persistent.
And that persistence is not fragile.

Outside, a car passes.

Inside, the diagram stabilizes. The wish for him to be visited by uncontrollable force fades into something cooler: Let him wrestle with his own reflections. That will be sufficient. She continues writing.

Ernest in his damp Airbnb imagining pursuit. Sylvia in her immaculate Holland Park flat imagining retribution. Both visited by forces entirely of their own making.

\section*{Constraints and Counterfactuals}

\lettrine{W}e are in a rather sullen pub near Angel in London. Late evening. Windows refract a road awash with careering cars outside going nowhere fast. \textit{Ernest Olber} leans against a cluttered table, wine glass in hand. \textit{Willard Jones} sits nearby, flipping through a printed article. The two men move easily within each other’s orbit: familiar, unhurried, the way seasoned conversationalists tend to be. Both have take a sabbatical and London suits their mood.

For Ernest, cosmology had transformed beyond explanatory models and measurements. He had rather come to see Einstein's field equations not as the geometrization of a freely falling force concept but as a framing prerequisite that made the language of physical fields legible. To him, Einstein's equations were little more than descriptive grammatic rules: not providing an impetus for Cosmic cause rather than imposing an insistence of its minutiae of particles to cohere.

Willard, always the logician, was not entirely opposed to this view. But where Ernest saw metaphysical structure, Willard sought a rub. An admirer of elegance yes, but only when it cast new insights on a process or collection or objects. Their conversations were a circling for a deeper truth neither ever pin down. Tonight, though, something in Ernest’s tone hinted at a pain not as yet fully declared. And it was Willard who, as usual, picked up the scab of this wound: half in frippery, half in concern.

\begin{itemize}[leftmargin=2cm, label={}]

\item[\textbf{Ernest:}]  
You’re missing the point, Willard. Those equations of Einstein aren’t about gravity-as-force. Rather they’re integrability conditions—metalanguage constraints. You want to find structure? Start where structure talks about itself.

\item[\textbf{Willard:}]  (Half-smiling).
That’s a logician’s move. Invoking Tarski to reframe field equations as commentary. But they’re still equations, Ernest. They still claim to govern.

\item[\textbf{Ernest:}]  
No, they \textit{permit} governance. Like syntax permits meaning. The real dynamics: the Rarita–Schwinger stuff, the spinors live in the object language. The Einstein tensor ensures that such things can \textit{coherently happen}.

\item[\textbf{Willard:}]  
You're saying curvature is grammar? That general relativity is the silent infrastructure? That’s... a bold downgrade.

\item[\textbf{Ernest:}]  
Not a downgrade. A shift. The RS equation needs a backdrop where the contact forms close. That’s the logic. Gravity ensures field equations aren't self-contradictory. It’s not describing matter rather it’s ensuring the description can even exist.

\item[\textbf{Willard:}]  
Fine. But it sounds like you're saying dynamics don’t matter. That the Einstein field equations are just bureaucratic rubber-stamps on deeper activity. 

\item[\textbf{Ernest:}]
Not rubber-stamps. Constitutional checks. You mistake the law for the courtroom drama. I’m talking about what makes the trial possible.

\item[\textbf{Willard:}]  
But you do this every time. You retreat to hierarchy. To formalisms that elevate structure above presence. 

\item[\textbf{Ernest:}]  
Because structure survives presence. Equations outlive collisions. Even you accept Bianchi identities as merely kinematical—fine. Then admit it: gravity’s not what moves us. It’s what makes motion describable. 

\item[\textbf{Willard:}]  
(sighs, glancing at the paper in his hand)  
All right, fine. Maybe I’ll concede that point.  
But don’t expect me to call your cosmological metaphors "grammar" in a referee report.  
There’s a limit to how much poetry one can smuggle into physics before it starts leaking meaning.

\end{itemize}

Ernest sets his glass down. The mood shifts. He picks up a copy of his \href{https://www.researchgate.net/publication/392402161_From_Spinor_Dynamics_to_Gravitational_Constraint_A_Tarskian_Interpretation_of_the_Einstein_Equations}{paper}, marked with citations.

\begin{itemize}[leftmargin=2cm, label={}]
\item[\textbf{Ernest:}]  
Wait No I won't accept that. This isn’t retreat. It’s a confession.  
Because if the Einstein equations are constraints: the metalanguage, then they don’t move things forward. They just prevent collapse.  
And if you think I’m talking about field theory alone, you’ve missed what this paper actually cost me.
\end{itemize}

Willard looks up as his eye had drifted away in the general direction of his mind. 

\begin{itemize}[leftmargin=2cm, label={}]

\item[\textbf{Willard:}]  
You mean... you can't mean the crash?

\item[\textbf{Ernest:}]  
What? Really—what? I meant all the constraints that came from it.  
(pause)  
The ones that shape you without guiding. Like guardrails that give way at the worst curve.
 

\item[]Silence. Willard shifts, unsure.

\item[\textbf{Willard:}]  
You say that as though it redeems you.  
As though becoming static made you less culpable.

\item[\textbf{Ernest:}] (flatly)
Not redeemed. Just resolved. 

\end{itemize}

\lettrine{T}{hey} proceed to revisit their prior orbit, sitting in a silence like a well-tempered clavier: measured, restrained, anything but resolved.  
A distant horn, perhaps another crash, echoes: a footnote to their communal failures.  
Willard, unaware of the full reasons behind Ernest's rising tetchiness, decides on a change of tack.  
He exhales, softly, and lets the weight of Ernest’s words settle. Then, as if flipping a page in his mind, he turns the paper in front of him and shifts from matters objective versus matters met to matters of implausibility.

\begin{itemize}[leftmargin=2cm, label={}]

\item[\textbf{Willard:}]  
Constraints are one thing. But surprise; that needs a measure.  
Otherwise, improbability is just poetry masquerading as principle. Meaningful probabilistic claims require a principled sampling scheme. Farey sequences provide such a scheme...” 
(Smiling) 

This is the core of it, Olber. It's not that the universe is improbably fine-tuned—it's that our framing lacks form. Without a defined space, there's no basis for surprise.

\item[\textbf{Ernest:}] (snorting) But I define awe, Willard. The universe isn’t improbable because of statistics—it’s improbable because it exists. So much form telenomically gathered together from so much apparent chaos.

\item[\textbf{Willard:}] But you’re relying on undefined counterfactuals. Sylvia—bless her razor clarity—shows that unless your “what ifs” live in a structured space, with metrics and form, you’re not doing physics. You’re indulging in metaphysics. We can’t claim improbability unless our “what-ifs” live in a structured geometry. Otherwise, we’re just playing dress-up with probability.
\end{itemize}
\begin{itemize}[leftmargin=2cm, label={}]

\item[]Willard taps the \href{https://www.researchgate.net/publication/392402161_From_Spinor_Dynamics_to_Gravitational_Constraint_A_Tarskian_Interpretation_of_the_Einstein_Equations}{paper}. Ernest watches, jaw tight.

\item[\textbf{Ernest:}] (quietly, the edge of his left eye - the faintest of twitches) Quoting Sylvia now?

\item[\textbf{Willard:}] Naturally. Her latest work. You've read it?

\item[\textbf{Ernest :}] I've read it before your footnotes began cluttered it up.

\item[] A pause. Ernest feeling an inner tension building.

\item[\textbf{Ernest:}] Do you even know how she came to all this hyperbolic number theory geometric stuff? To measures, conserved angles and rhyming triplets in homogenous co-ordinates?

\item[\textbf{Willard:}] I assumed that was just her thing. She won the Cramer Medal, didn’t she?

\item[\textbf{Ernest:}]The trouble with the laws of physics isn’t that they lie. It’s that they refuse to. When you're skidding toward collapse, they won’t invent another branch—another outcome. They stay cruelly true.

\item[\textbf{Ernest:}]  (raising a finger, theatrically):

\begin{center}
    
\textit{“From Farey lines the slopes align, in rational ascent,} \\
\textit{They pierce the sphere or kiss the curve, on how the cosmos bent—} \\
\textit{Triplets born where gradients meet, in measured firmament.”}

\end{center}

\item[] Willard lowers his eyes, unimpressed.

\item[\textbf{Willard:}] (dryly) 
You know, if you put half as much rigour into your metrics as your meter, the cosmos might just tune itself.

\item[\textbf{Ernest :}] (shrugging)  
Call it lyrical inference. The kind the field could use more of.
\end{itemize}

Ernest, mood souring stops swirling his wine as the pub begins to turn counter to it: presumably according to diktat of conserving the Universe's angular momentum. He s not removed from this isotropic madness. 
\begin{itemize}[leftmargin=2cm, label={}]

\item[\textbf{Ernest:}] (trying to be obtuse) It started with me. I broke something. Fifteen years ago. Coastal road. It was raining. I was arguing with her about Dirichlet domains or something. Anyway we were drunk anyhow, just drunk. We crashed. I crashed ...you know I crashed the car.  Pause

But. She was pregnant.
She was the one, the carrier of motion. She was never to come to term and I have never come to terms with that that.
\end{itemize}

Willard’s face visibly draws.

\begin{itemize}[leftmargin=2cm, label={}]

\item[\textbf{Ernest:}]  
Rehab was her returning to maths, for survival.  
She began calculating again in little fragmentary thoughts.  
In fairness though
she is not what she was.  
But even now, she outshines most of us.  
Back then... she was a constellation burning at full tilt.

\item[\textbf{Willard :}] (gently) That crash is the reason Sylvia is in a wheelchair? How come... I mean ..how did I not know this. You and me we talk abut more than just maths. Don’t we?  How did we never...

\item[\textbf{Ernest:}]  
We never found the moment. Or I avoided it. You, me—we’ve always been more about precessions than presence.

\item[\textbf{Willard:}] (the very personification of the unspoken as it happens) 
I thought I was better than all that. Than all the unspoken. 

\item[\textbf{Ernest:}]  
You are. You were. To be fair I thought we had more than that going for us.

\item[\textbf{Willard:}]  
So her preoccupation with projective geometry now: it’s a kind of reckoning with that moment?

\item[\textbf{Ernest:}]  
Yes. Her newer fascination with randomness—perhaps even with me, comes from that one point of collision: where my cosmology and her rational geometry met.  
% I was her hidden variable, as it were.  
% And, as it turns out, there were others neither of us had accounted for.  
So when I say the universe feels improbable, you call it poor framing.  
But I see her.  
I see the cost of its imbalance.

\item[\textbf{Willard:}] (softly)
Maybe both are true.  
The measure—of you as a man, of her... and of the memory you share.

\item[\textbf{Ernest:}]  
Some measures,  
the ones that truly matter,   
we only understand in hindsight.
\end{itemize}

Ernest pauses. The glass in his hand tilts slightly. Not for affectation more just the gravity of it all catching up with him.

\begin{itemize}[leftmargin=2cm, label={}]

\item[\textbf{Ernest:}]  
(Now barely audible). She couldn’t have children after.  That detail never made it into the medical notes. But I remember the consultant’s phrasing—“structural trauma inconsistent with future viability.”  

\end{itemize}

Ernest swirls the wine, doesn’t drink. Glaring now present-mindedly at the Newton bucket experiment unfolding in his grasp.

\begin{figure}[H]
\centering
\includegraphics[width=0.7\linewidth]{Illustrations/vign-ErnestGlass.png}
\caption{Ernest all minded }
\end{figure}


\begin{itemize}[leftmargin=2cm, label={}]
\item[\textbf{Ernest:}]  
And I couldn’t either.  
Not because I was hurt.  
Because I never found the language to separate the want from the guilt.

\item[\textbf{Willard:}]  
You mean you chose not to?

\item[\textbf{Ernest:}]  
I don’t know. What’s choice when you’re carrying with you the denominator of someone else’s life?

\end{itemize}
A long silence passes that was best not be filled.

\begin{itemize}[leftmargin=2cm, label={}]

\item[\textbf{Ernest:}]  
I wasn’t meant to be a father. But I was—once. In that moment.  And then I wasn’t. And neither was she.

\end{itemize}
\section*{The pain of our penumbras}

\lettrine{E}dward opens his mouth to speak, but merely forms an Oh. The conversation clots as they both lock onto a set of elevated voices in the room of which they still reside. Mutually aware, a relative stasis returns: all clauses briefly arrested, all trespass momentarily absolved.
 Willard nods; more to the space between them than to Ernest directly. He knows he could say more. Knows, too, that whatever he says might feel wrong in the morning. He gathers Sylvia’s paper all creases, margins inked and walks out into the cool hub-hub of the high street to go catch his bus, leaving the room and Willard to ruminate on his broken silence.

Willard happens to catch a ride back to his temporary digs on the \emph{43} bus, half-unaware, he grabs a seat just behind the driver, \emph{entitled}, as he sometimes felt, to the most accessible of seats. Age had earned him that much, he reckoned on, even if no badge or stick pronounced it. A couple of stops in the journey a grissly \emph{old fella} clambers aboard, stooped like a penguin, folded over like a letter C eyes flicking over the bus corridor, searching for the release of a temporary repose.  Slumping he barks, loud enough for the whole bus to hear:
\begin{quote}
\emph{``Is this hat anyone's?''}
\end{quote}

Willard looked down at the woolly beanie lying on the chair in front of him.
\begin{quote}
\emph{``No,''} gently,\emph{``it's been there a while, so I reckon it's yours now.''}
\end{quote}

The grizzly grumbled and muttered, jamming the cap on his sweaty bonce with the deftness which lent credence to the authenticity of his later declaration of being a rather  handy pick-pocket-er . After a moment he struck up a lighter, its flame trembling in the unsteady cup of his hand, and drew on a cigarette rolly. Its acrid smoke swelling through the bus, a welcome masking of the formic odour of the accumulated prostate leaching between his thighs. Willard leaning towards him, speaking softly, spontaneously - unaware of London's non interference etiquette:
\begin{quote}
\emph{``That's not fair to the driver—you’re putting him in a difficult position—you know you shouldn't be smoking on a bus.''}
\end{quote}
The reply was immediate and as bitter as any city dweller would attest :
\begin{quote}
\emph{``Don't you be telling me what to do, mate. Fuck yourself off.''}
\end{quote}

Heads briefly lifted across the bus. Willard's chest missing a beat he was stilled.

\begin{quote}
\emph{``I used to be in the SAS—I’ve killed fifty terrorists you know —I don’t mind doing in another.''}
\end{quote}

He proceeded to ramble, his voice then sinking into a ragged murmur, then in a moment of audience awareness catching himself for a moments:
\emph{``Me. I’ve fallen on bad times.''}

Then continues, airing his dirty laundry to the ten or so unwilling inhabitants of who are variously find fascination in window refractions and reaching for ear pods. No one speaks as is the way these days. Everyone receiving his pain within their breathless breasts.

Willard taking the lead from the natives, himself now stares out at the rain beginning to streak the glass, refracting the neon lights of breaded chicken shop derivatives. Their polyunsaturated sheen belying the horror of the daily massacre that produced protein in a bap. A dull ache now forms behind his ribs: not of panging, not of fear, but slow recognition.

\begin{quote}
\emph{At our age we were all someone once. We are all half-castes. Penumbric shadows of our former radiant selves.}
\end{quote}

He reflects how he is in the latter period of his persistence. That when we subsist uncomfortably in the penumbra \emph{of our former glorious self}, clinging to the less shadowy umbra we once cast, desperate to declare to others our former glories before the world never realizes we were ever here.

The bus sighs to a halt. The old man shuffles off drizzling into the drizzle, his no little paucity of a pall of rank hue trailing behind him. Willard watches him reach for the nearest bin at the bus stop as his image strokes left into the dark and whispering to himself begrudgingly: \emph{``So am I.''}

Willard now, more than a little overcome, steps off two stops early, walking slowly, the rain with impatient fingers tapping at his collar. He turns towards a railway bridge, stroking its long steel arc thinking of what lies beneath: an ordained track of time-like trajectories. That for Sylvia, that for us  massive objects with built obsolescence:
\begin{quote}
\emph{``Nothing with mass persists in form. Our inertia casts impermanent shadows, and with our radiative decay, those shadows dilute to penumbra with age.''}
\end{quote}

Willard was picturing her little hands on the wheels of her chair, gilding her motion with patient precision—every turn being earned. He had not been the one behind the wheel the night everything changed, but he felt as though he was in some little way. Her body broke that night, but her spirit annealed; rebuilt with steel in her jaw while he drifted like silt.

She had every right to hate Ernest, and by vicarious association him, yet she never appeared to. The cruelty that hid in her solemnity of quiet suffering.

He leaned against the cold railing and let the wind sting his eyes. The old man's words echoed:
\begin{quote}
\emph{``I’ve also fallen on bad times.''}
\end{quote}

Willard whispered into the rain:
\begin{quote}
\emph{``You were someone once too, Sylvia.''}
\end{quote}

% \newpage

Willard eventually returns to his crappy apartment on the Holloway road. The lamp is still on. Sylvia’s paper lies open where he gently placed it to a page marked with dense annotations. He is a figure of angular distributions sitting beside a wine glass.  Breaking symmetry he reposes a phone to his ear, clearing throat in readiness for the softest of his tones.

\begin{itemize}[leftmargin=2cm, label={}]

\item[\textbf{Willard:}]  
Sylvia? It's Willard. I hope I’m not waking you.

\item[] A pause. A tired but alert voice answers.

\item[\textbf{Sylvia:}]  
Not at all. Are you calling about the paper?

\item[\textbf{Willard:}]  
Yes.  
I should’ve read it properly when you sent it. I skimmed, too quickly. But tonight...  
Tonight I couldn’t stop. Especially the Farey sequence sections. Figure 3 really held me.

\item[] A slight rustle on the line. She’s listening.

\item[\textbf{Willard:}]  
You’re doing more than analysis. That convergence pattern—the way the Farey slopes stabilize distribution—it’s like watching geometry instruct measure.  
You say it best: \textit{"Probability depends on measure."}  
This time, that line landed differently.

\item[] Another pause. He flips a page.

\item[\textbf{Willard:}]  
And the cross-ratios: those centroidal rays through the triangles.  
The contrast between affine and conformal mappings and how one preserves predictive power, the other flattens it. It’s elegant work, Sylvia. Beautifully argued.

\item[\textbf{Sylvia:}]  
Thank you Willard.  
I wasn’t sure anyone would follow the Delaunay part. Or the parity-null results as they were a bit off pis

\item[\textbf{Willard :}] (delicately)
I did. I really did.  
(Pause)  
And I heard something else.  
Not in the maths.  
Between the lines.

\item[\textbf{Willard:}]  
You said it best: \textit{“Probability depends on measure.”}  
But if the Farey gradients define the space...  
Then geometry is instructing the measure—  
And by transitivity, it’s instructing probability.  

\item[] A pause. Sylvia arches an eyebrow.

\item[\textbf{Sylvia:}]  
Careful, Willard. You’re starting to sound like Ernest.

\item[\textbf{Willard :}]  (half-laughing)
Am I? God help me.  
Next I’ll be quoting Tarski and calling tensors 'metalanguage.'

\item[\textbf{Sylvia :}]  (dryly)
You’ve spoken to him recently, haven’t you?

\item[\textbf{Willard :}] (evasively) 
Define “recently.”


\item[] Silence on the line. Long enough for him to think she might hang up.

\item[\textbf{Sylvia :}] (measured) 
Ask no question, I will tell you no lies.

\item[\textbf{Willard :}] (quietly) 
Fair. Especially this late.

\item[] A pause. The line hums with everything unsaid.

\item[\textbf{Sylvia :}] (softer now)  
But if you called to say something, Willard...  
Then say it. Even if you don’t know what it is yet.

\item[] He exhales, a sound like a page turning.

\item[\textbf{Willard:}]  
I read your paper, Sylvia.  
But maybe... I was reading for something else.

% \item[] Silence stretches across the distance.

\item[\textbf{Sylvia :}]  (gently)
Then read again.  
And this time—read as yourself.

\item[\textbf{Sylvia :}]  (wearily)
What is it you really want to say, Willard?

\item[\textbf{Willard :}]  
(softly) Ernest told me. About the drive. About the consequences.  
I didn’t know. I didn’t even know to be sorry. But I am. Deeply.

\item[\textbf{Sylvia:}]  
Is that what you really called to say?

\item[\textbf{Willard:}]  
He said you started calculating again in rehab.  
I don’t know what to say, Sylvia—except that your work... it may have taken time to return, but it’s not just precise. It’s brave.

\item[\textbf{Sylvia :}]  
(flatly) Ernest said?  
And here I thought your compliments were borne of admiration, not pity.

\item[\textbf{Willard:}]  
No—no. They stand.  
I kept looking at your Farey-based triplet-generation procedures... and all I could think of was how beautifully ordered and rigorous they are. How measured. Ant those transofrmation to and from projective space:  
like someone trying to impose a Mobius structure on something that wouldn’t stop spinning.

\item[\textbf{Sylvia:}]  (quiet)
I remember the spin Willard. Not the impact though.

\item[\textbf{Willard :}]  
(gently) Right. Yes.  
Sylvia really. I didn’t know about the baby.  
I’m sorry. Not just for reading late—for not seeing sooner.

\item[] A breath. Almost words. But not yet.

\item[\textbf{Willard:}]  
Your paper isn’t just about geometry, anyway Sylvia.  
It’s about choosing spaces where things can be defined.  
Where probabilities aren’t just guesses. Reaching.  
And maybe... grief works like that too.

\item[] She says nothing. He lowers his voice.

\item[\textbf{Willard:}]  
I’ll cite it with care Sylvia.  
But I’ll read it again—and again.  
Not just because it’s sharp. But because it came from the place beyond error bounds.

\item[\textbf{Sylvia :}] (faintly) 
It’s still just mathematics.

%%%%%%
\item[\textbf{Willard:}]  
Anyway—what can a lot of words say  
that such information-rich mathematics doesn’t already say better?  
More sharply. More cleanly. I should just be... watching you reveal those truths  in \LaTeX character form.
(grimacing to himself and audibly wincing)
God, that sounded creepier than I meant it.  
I mean—watching the ideas unfold. Not—  
not watching \textit{you}, watching.  
You know what I mean.

% \item[] Silence. Then—just possibly—a dry, amused exhale from Sylvia.


%%%%%%%%%
\item[\textbf{Sylvia:}]  
Relax, Willard. I’ve had creepier admiration from referees.  
At least you read past the abstract.
And let me guess...  
You talked to Ernest recently?
%%%%%%%%%%%
\item[] Silence. Then a breath—longer than before.

\item[\textbf{Sylvia:}]  (softly)
You’re the first to call me about it.  
Not email. Not footnotes.  
A voice. Yours.

\item[] Willard swallows, unsure whether to speak.

\item[]  (continuing, barely audible)
I thought maybe I’d gone quiet in the field.  
Too long out of step.  
Too long... being the one people don’t look at.

\item[] A pause hangs, close and exposed.

\item[\textbf{Willard:}]  
Sylvia, really your paper isn’t just rigorous,  
it’s luminous: like a map that remembers more than one path.

\item[] A quiet exhale. Almost a sob, swallowed down.

\item[\textbf{Sylvia:}]  
You have no idea how much I needed to hear that.  
Not because I doubted my analysis.  
But because I wasn’t sure... if anyone still saw me as part of the field.

\item[] Willard doesn’t answer immediately. He lets her be heard.

\item[\textbf{Willard:}]  
I see you Sylvia.  
More clearly than ever.

%%%%%%%%%%%%%
\item[] A breath, almost a thank you—but caught.

\item[\textbf{Sylvia :}] Well... that’s kind of you, Willard. (more composed)
It’s late. You should get some sleep.  
And I have... annotations to finish.

\item[] The warmth doesn’t vanish—but it’s sealed, respectfully.

\item[\textbf{Willard:}]  
Of course.  
Goodnight, Sylvia.

\item[] A beat. No reply.  
Then, just before the click:

\item[\textbf{Sylvia :}]  (very faintly)
Goodnight, Willard.
%%%%%%%%%%%%%%%%%

\end{itemize}

The call ends. The screen dims. Willard remains seated, phone still in hand, as if listening for something more. 

Just sitting there.  Not reading. Not calculating.  
Just thinking about all that he never says.

\begin{itemize}[leftmargin=2cm, label={}]

\item[\textbf{Willard :}]  (in a whisper to himself)
Who did I ever invite to stand still in my secret garden with me...  
just long enough to notice something growing or retreating into dormancy? Not Ernest. Not really.  
Not in all those years.
\end{itemize}
He presses his thumb against a corner of Sylvia’s paper, not turning the page yet. He flips it.  Then flips it back. Too quick either time to have absorbed anything new. A self-conscious rhythm.  A dance performed for an audience long gone, or rather never there. 

Eventually he does read it again. Slowly, more laboured this time. Letting each Farey slope mark time on his own rhythms of regret.



\section*{Rational Geometric Probability Structures}

Farey sequences define dense, statistically regular sets of rational slopes on the interval \( [0,1] \). As the Farey order \( N \to \infty \), these slopes approximate a uniform angular distribution on the projective line, making them a principled foundation for measure-theoretic constructions in geometric probability.

\begin{figure}[H]
  \centering
  \includegraphics[width=0.45\textwidth]{Illustrations/comparisonFareyArithmeticalRationalSlopes.png}
  \caption{Farey sequence (top) vs. arithmetic slopes (bottom) — order and bias contrasted.}
\end{figure}

% >>>>>> ANNEX B (cut-sheet v2): Farey-sequence measures and the cross-ratio invariant — block commented out below, pointer left in text <<<<<<
\textit{(Sylvia's Farey-sequence measure on Pythagorean triples, and the cross-ratio it rests on, are set out in full in Annexe B.)}

% \subsection*{Farey Sequences as Rational Measures}
% By projecting lines of rational slope from a fixed point (e.g., \( (-1, 0) \)) on the unit circle, Farey sequences generate rational intersection points, corresponding to solutions of \( x^2 + y^2 = 1 \). Clearing denominators yields Pythagorean triples:
% \[
% a = q^2 - p^2, \quad b = 2pq, \quad c = q^2 + p^2,
% \]
% where \( m = \frac{p}{q} \in \mathcal{F}_N \), \( \gcd(p,q) = 1 \), and \( p \not\equiv q \pmod{2} \). This arithmetic structure ensures that the resulting triples are both systematically generated and statistically significant.
% 
% \begin{figure}[H]
%   \centering
%   \includegraphics[width=0.45\textwidth]{Illustrations/fareyCircle.png}
%   \caption{Primitive (blue) and non-primitive (red) \href{https://colab.research.google.com/drive/1U_w19FQYhnVYBlGdL9P90rSszdAv0SWW?usp=sharing}{Pythagorean triples}}
% \end{figure}
% 
% \subsection*{Cross Ratios and Projective Invariants}
% Farey slopes interpreted as angles define rays whose angular cross ratios tend to cluster around \( \frac{4}{3} \), a hallmark of projective harmonic configurations. Unlike arithmetic slope sequences, which exhibit erratic behavior, Farey-based cross ratios suggest an emergent regularity—one that supports the construction of projectively invariant measures on rational geometric configurations.
% 
% \begin{figure}[H]
%   \centering
%   \includegraphics[width=0.45\textwidth]{Illustrations/sweepFarey.png}
%   \caption{Convergence of primitive triple probability from Farey slopes toward \( \frac{3}{\pi^2} \).}
% \end{figure}
% 
% 
% \subsection*{The Cross-Ratio: A Projective Invariant}
% 
% One of the central objects in projective geometry is the \emph{cross ratio}, that encodes how four collinear points (or concurrent rays) relate to one another in a way preserved by Möbius transformations. {\textit{Cross ratio}}\label{CrossRatio} of four\footnote{With three points \(A, B, C\), any linear transformation preserves their order and relative spacing. With four points \(A, B, C, D\), the relationship becomes richer: their cross-ratio captures a projective invariant.} distinct collinear points \(A, B, C, D\) defined as:
% \[
% \xi(A, B, C, D) = \frac{(A - B)(C - D)}{(B - C)(D - A)}.
% \]
% remains unchanged under a projective projection. In the figure below, points \(A, B, C, D\) lie on the real line, while rays from a common point (the “eye”) project them to new points \(A', B', C', D'\) on a diagonal line. Though the coordinates differ, the cross ratio is preserved.
% 
% \begin{figure}[H]
% \centering
% \includegraphics[width=0.5\textwidth]{Illustrations/emanate.png}
% \caption{The projection of four points onto a new line preserves their cross-ratio.}
% \label{fig:projCross}
% \end{figure}
% 
% \begin{minipage}{.48\textwidth}
% \centering
% \textbf{Original Points (X, Y)}
% \begin{tabular}{|c|c|c|}
% \hline
% Point & X & Y \\
% \hline
% A & 2.00 & 0.00 \\
% B & 3.00 & 0.00 \\
% C & 5.00 & 0.00 \\
% D & 8.00 & 0.00 \\
% \hline
% \end{tabular}
% \end{minipage}%
% \begin{minipage}{.48\textwidth}
% \centering
% \textbf{Projected Points (X, Y)}
% \begin{tabular}{|c|c|c|}
% \hline
% Point & X & Y \\
% \hline
% A' & 1.71 & 1.17 \\
% B' & 2.39 & 1.63 \\
% C' & 3.50 & 2.40 \\
% D' & 4.75 & 3.25 \\
% \hline
% \end{tabular}
% \end{minipage}
% 
% \smallskip
% 
% Both sets yield the same cross-ratio:
% \[
% \xi = \frac{(2 - 3)(5 - 8)}{(3 - 5)(8 - 2)} = \frac{(-1)(-3)}{(-2)(6)} = -\frac{1}{4}.
% \]
% 
% \noindent This invariance under projection is foundational to why the cross-ratio is used to define harmonic division. 
% 
% \textbf{Farey Case:} When slopes from Farey sequences are interpreted as angles of rays from a fixed base point, the angular cross ratios formed by quadruples of directions tend to cluster around \( \frac{4}{3} \). This value is associated with harmonic division and suggests a deeper regularity in rational constructions.
% 
% \begin{figure}[H]
% \centering
% \includegraphics[width=0.55\textwidth]{Illustrations/sweepFarey.png}
% \caption{Clustering behavior of cross-ratios from Farey slopes.}
% \end{figure}
% 
% Together, these observations hint that projective invariants like the cross ratio are not merely geometric curiosities but bridges between algebraic sequences and geometric structure—appearing as "fixed points" in the space of rational configurations.
% 
% Define now \( F_n \) as the Fibonacci numbers and consider four consecutive values \( F_{n-3}, F_{n-2}, F_{n-1}, F_n \) whose cross-ratio is:
% \[
% \xi_n = \frac{(F_{n-3} - F_{n-2})(F_{n-1} - F_n)}{(F_{n-2} - F_{n-1})(F_n - F_{n-3})}.
% \]
% as \( n \to \infty \). Binet's formula for golden ratio \( \phi = \frac{1 + \sqrt{5}}{2} \) reads
% \[
% F_n = \frac{\phi^n - \hat{\phi}^n}{\sqrt{5}} \quad \text{where } \hat{\phi} = \frac{1 - \sqrt{5}}{2}.
% \]
% Taking the limit as \( n \to \infty \), we have \(
% \lim_{n \to \infty} F_n \approx \frac{\phi^n}{\sqrt{5}}.
% \) since \( |\hat{\phi}| < 1 \) and \( \hat{\phi}^n \to 0 \) and our cross ratio thus reads as:
% \[
% \lim_{n \to \infty}\xi_n = \frac{(\phi^{n-3} - \phi^{n-2})(\phi^{n-1} - \phi^n)}{(\phi^{n-2} - \phi^{n-1})(\phi^n - \phi^{n-3})},
% \] This in turn satisfies the fundamental identity:
% \[
% \phi^2 = \phi + 1,
% \]
% implying a recursive structure:
% \[
% \phi^2 = \phi^1 + \phi^0, \quad \phi^3 = \phi^2 + \phi^1, \quad \phi^4 = \phi^3 + \phi^2, \quad \ldots \phi^n = \phi^{n-1} + \phi^{n-2}
% \]
% \noindent
% We also have \(\phi^3 - 1 = 2\phi\) since:
% \[
% \phi^{n} = \phi^{n-3} \cdot \phi^3\quad 
% \phi^{n-1} = \phi^{n-3} \cdot \phi^2, \quad 
% \phi^{n-2} = \phi^{n-3} \cdot \phi,
% \]
% and thus:
% \[
% \phi^3 = \phi^1 \cdot \phi^2 = \phi(\phi + 1) = \phi^2 + \phi = 2\phi + 1.
% \]
% \noindent
% Applying the identity \( \phi^{n-k} = \phi^{n-3} \cdot \phi^{k} \) for \( k = 0, 1, 2, 3 \), the invariant reads:
% \[
% \begin{aligned}
% \lim_{n \to \infty}\xi_n &= \frac{(\phi^{n-3} - \phi^{n-3} \cdot \phi)(\phi^{n-3} \cdot \phi^2 - \phi^{n-3} \cdot \phi^3)}{(\phi^{n-3} \cdot \phi - \phi^{n-3} \cdot \phi^2)(\phi^{n-3} \cdot \phi^3 - \phi^{n-3})} \\
% &= \frac{\phi^{n-3}(1 - \phi) \cdot \phi^{n-3} \cdot \phi^2 (1 - \phi)}{\phi^{n-3} \cdot \phi (1 - \phi) \cdot \phi^{n-3} (\phi^3 - 1)} \\
% &= \frac{\phi^{2(n-3)} (1 - \phi)^2 \phi^2}{\phi^{2(n-3)} (1 - \phi) \phi (\phi^3 - 1)}.
% \end{aligned}
% \]
% Cancelling common terms we are left with:
% \[
% \frac{(1 - \phi)\phi^2}{\phi(\phi^3 - 1)}= \frac{(1 - \phi)\phi^2}{\phi \cdot 2\phi} = \frac{(1 - \phi)\phi^2}{2\phi^2} = \frac{1 - \phi}{2}= \frac{2 - (1 + \sqrt{5})}{4}.
% \]
% 
% \medskip
% 
% \noindent
% \textbf{Final Result:}
% \[
% \boxed{
% \lim_{n \to \infty} \xi(F_{n-3}, F_{n-2}, F_{n-1}, F_n) = \frac{1 - \phi}{2}  = \frac{1 - \sqrt{5}}{4}=\frac{\hat{\phi}}{2}
% }
% \]
% 
% 
% 
% 
% 
% \newpage
% >>>>>> end of annexed block <<<<<<
\newpage
\lettrine{J}{ohn} O’Leary had always been a mathematical wanderer with his surname evoking, with uncanny aptness, both Leray’s sheaf theory and a tendency to leer at the work of others. From Leray to Grothendieck, he traced conceptual paths with a zeal that often outpaced his originality. His eye, ever roving, did not discriminate between beautiful ideas and vulnerable colleagues. For those in his field, particularly Sylvia, his presence in one’s orbit felt less like a fellow traveler, more like rogue debris accreting, absorbing, and repackaging. Sylvia had been developing her Farey sequence measure theory of Pythagorean triples, a delicate, exacting little framework. But Leary, ever alert to novel terrain was looking to retrofit with the heavy machinery of line bundles and Hopf fibrations, could see an opportunity. And with his usual opportunism, he set out to sheaf-ify it, by lifting it from the familiar intersections of rational lines on the unit circle to the more exotic terrain of non-Euclidean Riemannian spaces.
% \begin{figure}[H]
% \includegraphics[width=1.0\linewidth]{Illustrations/vign-john.png}
% \caption{John O'Leary reading with some intensity }
% \end{figure}

% It was the whispers of Eleanor’s gathering at the Merriweather Library that had drawn him in, his curiosity piqued by the potential connections between classical geometry, sheaves, and the mysteries of intersecting lines and circles.

% \subsection*{Pencils of Rational Lines and Germs of Triples}
% Leary’s fascination began with a simple setup: a pencil of rational lines emanating from (-1,0) in the plane, intersecting the unit circle \(x^2 + y^2 = 1\). Each line had a slope \(m\), given by the equation:
% \[
% y = mx + m.
% \]
% The intersections of these lines with the circle produced points that, when rational, corresponded to Pythagorean triples:
% \[
% x^2 + (mx)^2 = 1 \implies x = \frac{1}{1+m^2}, \quad y = \frac{m}{1+m^2}.
% \]
% For rational \(m = \frac{p}{q}\), the coordinates \((x, y)\) scaled to integers, generating triples such as \((3, 4, 5)\).



 \begin{figure}[H]
\includegraphics[width=1.0\linewidth]{Illustrations/vign-arnieFish.png}
\caption{John learing at a fish}
\end{figure}


What, Leary mused, would happen if the flat Euclidean canvas that underpinned Sylvia’s pencil of concurrent lines were generalized? Suppose the foundational space were no longer flat but curved—Riemannian. What if, instead of the unit circle, the intersections took place on a saddle-shaped surface (hyperbolic space), or across the bulge of a sphere (elliptic geometry)? How would the structure of intersections transform? \footnote{Topological Note: The classical Jordan Curve Theorem asserts that any simple closed curve in the plane divides it into an interior and exterior region. In our construction, the sub-circle lies not in the Euclidean plane but on the 2-sphere. Fortunately, a generalized version of the Jordan Curve Theorem applies: a simple closed curve on any two-dimensional manifold  still separates the surface into two connected regions. Thus, although our "unit circle" is not a great circle and lacks a planar embedding, it retains an interior/exterior structure on the sphere. This topological property underwrites our ability to speak meaningfully of geodesic intersections, enclosures, and angular deviation in this curved setting.}
 Could Sylvia’s sheaf of rational sections survive the journey into curvature?

In sheaf-theoretic language, the pencil of lines becomes a section of a line bundle over the unit circle. On the Euclidean plane, the sheaf's \emph{germ}—its local behavior—was straightforward, reflecting rational parameterizations and trivial twisting. But curvature introduced complexity. The bundle could no longer remain flat. On a hyperbolic surface, for instance, lines intersect not circles but geodesic boundaries. These intersections are no longer governed by rational ratios, but by hyperbolic trigonometric structures,\(\sinh\) and \(\cosh\). The sheaf, adapting to its new terrain, would encode these changes in its transition functions, incorporating curvature directly in terms like:
\[
\theta^a \wedge \theta^b R_{ab},
\]
where \(\theta^a\) denote hyperbolic analogs of the basis one-forms, and \(R_{ab}\) captures the curvature’s effect on local behavior.

The germs of the sheaf in hyperbolic space would encode a dispersive rather than algebraic character. Lines, no longer rational, would chart geodesic flows. The Pythagorean triples, once a disciplined parade of integers, would now dance to the tune of negative curvature, reinterpreted as hyperbolic echoes of their former selves.

\begin{figure}[H]
\centering
\includegraphics[width=1.0\linewidth]{Illustrations/vign-learyPickles.png}
\caption{John O'Leary staring down a pickle}
\end{figure}

Leary stood abreast his prized collection of pickled peppers; a gleaming alignment of jars, a chromatic suggestion: could the entanglement of flavors in a well-brined pepper mirror the mathematical twisting of line bundles? Just as pickling transformed a plain pepper into something layered and unexpected, so too did twisting in sheaf theory reshape flat, rational structures into richly curved geometries. It wasn’t such a stretch. Curvature could transmute Sylvia’s clean intersections with the twisting of a sheaf over a base space altering local-global relationships. 

\paragraph{Hopf Bundles and Twisted Rationality.}

Leary’s thoughts turned to the Hopf bundle,
\[
S^3 \to S^2,
\]
a \textbf{fibration} with fiber \( S^1 \), mapping the 3-sphere onto the 2-sphere via twisting circles. This shift—from untwisted Euclidean fibers to nontrivially entangled circles offering up a natural framework for understanding sheaves over curved spaces.

In this setting, Sylvia's rational slopes generalize to rational directions within a twisted line bundle and notions like \textbf{primitivity} and \textbf{cross-ratio} become curvature dependent invariants. The geometry is no longer globally flat, and the sheaf’s local data that are the germs, encode this global twist via nontrivial transition functions of the bundle.

Bahti was ambitiously drawing parallels to the Frobenioid framework,\footnote{Introduced by Mochizuki, Frobenioids generalize line bundles with Frobenius actions into categorical structures capturing arithmetic through morphisms. Bahti’s adaptation imagines rational geodesics encoding curvature-aware arithmetic within a structured sheaf.} where number, geometry, and category interact. Sylvia’s constructions, reframed, suggest a directionally-structured bundle of numbers—where geodesics inherit curvature and number-theoretic structure alike.

Over hyperbolic bases, fibers bend apart in negative space; over spheres, they close along great circles. In both cases, rational parameterizations give way to localized expressions shaped by curvature and topological entanglement.

% >>>>>> ANNEX B (cut-sheet v2): line bundles & Frobenoidal structures (Leary's machinery) — block commented out below, pointer left in text <<<<<<
\textit{(Leary's sheaf-lifted version of her framework --- line bundles and Frobenoidal structures --- is reproduced, with attribution restored, in Annexe B.)}

% \subsection*{Of Line Bundles and Frobenoidal Structures}
% 
% The classical sheaf of concurrent rational lines intersecting a unit circle is re imagined on curved surfaces by replacing lines with geodesics. On a sphere, these geodesics are great circle arcs, emanating from a base point on a sub-circle embedded along the equator. Each such arc corresponding to a Farey slope reframed as a tangent direction whose intersections with the sub-circle form a curved analogue to the rational structures seen in Sylvia's plane.
% 
% This geometry is governed not by linear equations, but by spherical trigonometric constraints: periodic functions like \(\sin\) and \(\cos\) dictating the behavior of intersections, and the resulting sheaf exhibiting finite cyclic symmetries. Its germs encoding local data reflective of spherical curvature. Pythagorean triples, once rational and flat, now emerge as tessellated spherical triangles.
% 
% Within this framework, sheaves cease to be flat and begin to resemble line bundles on curved spaces whose \textbf{transition functions} capture the topology of curvature itself. Such constructions are a nod toward Frobenioid structures,\footnote{See Mochizuki’s formulation of Frobenioids, where line bundles with Frobenius morphisms are reinterpreted categorically. Bahti adapts this perspective for rational geodesics under curvature, treating them as morphic data within a bundle-theoretic arithmetic.} where number theory, topology, and categorical structure converge.
% 
% \begin{figure}[h]
%   \centering
%   \includegraphics[width=0.8\textwidth]{Illustrations/geodesicIntersection.png}
%   \caption{Farey-generated \href{https://colab.research.google.com/drive/1l-5OZV4MtRtw3oZykz6-j3rlrOtXAWSB?usp=sharing}{geodesics} from a base point on the sphere. Curvature disrupts Euclidean intersection patterns, motivating a projective reinterpretation.}
% \end{figure}
% 
% In projective geometric\footnote{The emphasis on projective embeddings and invariant cross-ratios hints at a deeper categorical structure underlying Sylvia's triangle constructions. Just as certain generalizations in arithmetic geometry such as \emph{Frobenioids} recast valuations, divisors, and Frobenius actions in the language of morphisms and monoidal categories, so too can one view Bahti's maps from rational slope data into projective conic classes as morphisms between structured spaces. These embeddings preserve not only geometric shape but also arithmetic regularity, suggesting that a richer categorical account of triangle structure which is sensitive to parity, primitivity, and curvature, may be possible.} terms, this construction represents a move from Euclidean to spherical measures. Deviations from the plane-case intersections encode curvature-induced distortions.
% 
% Leary is wondering what Bahti has long since already pondered: What becomes of Pythagorean triples under constant-curvature deviation? How do rational parametrizations behave on such curved manifolds? Are these intersection points governed by spherical harmonics or bundled structures like those found in Hopf fibrations? 
% 
% \begin{figure}[H]
% \centering
% \includegraphics[width=0.5\linewidth]{Illustrations/vign-farmer.png}
% \caption{John O'Leary envisaging his father with sheaves of germs of wheat}
% \end{figure}
% 
% Leary closed his notebook. His thoughts spun with rational pencils, twisted sheaves, and the topology of the Hopf bundle. But his real insight was this: when the rigid Gaussian integer lattice was allowed to flex like a hammock sagging into curvature as the tidy rational structure begins to slip. The arithmetic was no longer strictly rational, but \emph{realised}, even gently irrational, shaped by the yielding of the underlying space.
% 
% “Perhaps,” he mused, “mathematics isn’t about rational order, but what unfolds when the substrate space of (be it of Gaussian or Eisenstein integers) itself yields beneath it.”
% 
% Plucking another pickled pepper from the jar he steps into the city’s amber light. As the afternoon wanes, with notebook under his arm he pictures his father holding a sheaf of wheat while he walks to the Café Equilibrium, a known refuge for unsettling ideas where he was to meet with Eleanor who had called an in person meeting.
% 
% 
% \newpage
% 
% >>>>>> end of annexed block <<<<<<
